The \code{mem\_perm} class represents
general memory page permission for the given memory page in the process.


\begin{apient}
Process::mem_perm::read Process::mem_perm::write Process::mem_perm::execute
\end{apient}

\apidesc{
The Process::mem\_perm::read, Process::mem\_perm::write, and
Process::mem\_perm::execute, just as their names imply, respectively represent
read, write, and execution permission of given memory page. 
}

\begin{apient}
mem_perm(bool r, bool w, bool x)
\end{apient}

\apidesc{
These constructors provide a convenient way to create the specific memory
permission for the given page.
}

\begin{apient}
bool getR() const
bool getW() const
bool getX() const
\end{apient}

\apidesc{
These functions return true if the given memory page has read, write, and
execution permission, respectively, and false otherwise.
}

\begin{apient}
bool isNone() const
bool isR()    const
bool isX()    const
bool isRW()   const
bool isRX()   const
bool isRWX()  const
\end{apient}

\apidesc{
These functions return true if the permission of given memory page is
NO\_ACCESS, READ\_ONLY, EXECUTE, READ\_WRITE, READ\_EXECUTE, and
READ\_WRITE\_EXECUTE, respectively, and false otherwise.
}

\begin{apient}
Process::mem_perm& setR()
Process::mem_perm& setW()
Process::mem_perm& setX()
\end{apient}

\apidesc{
These functions enable read, write, and execution permission for the given page,
respectively, and return this mem\_perm.
}

\begin{apient}
Process::mem_perm& clrR()
Process::mem_perm& clrW()
Process::mem_perm& clrX()
\end{apient}

\apidesc{
These functions disable read, write, and execution permission for the given
page, respectively, and return this mem\_perm.
}

\begin{apient}
bool operator<(const mem_perm& p) const
\end{apient}

\apidesc{
This function returns true if this mem\_perm is less than p according to the
notation that read permission encodes to 4, write, 2, and execute, 1, and false
otherwise.
}

\begin{apient}
std::string getPermName()
\end{apient}

\apidesc{
Return the memory permission name for this mem\_perm.
}